% GENERATED FILE. Edit canonical lessons, then run npm run generate:books.
% canonical-source-hash: fnv1a64:a9f8fa09f79b6df8
% canonical-lessons: SA-C310-ghatana, SA-C310-apatati, SA-C310-vivaranam, SA-C310-prakasayati, SA-C310-visesah

\chapter{What Happened Today}
\label{ch:sa-what-happened-today}
\hlchaptermodality{310}

\begin{chapteropening}
\textbf{By the end of this chapter:} \emph{I can say what happened and give the details.}

Complete the last lesson of chapter 310: I can say what happened and give the details.
\end{chapteropening}

\section[\texorpdfstring{\sk{घटना} (ghaṭanā)}{घटना (ghaṭanā)}]{\sk{घटना} (ghaṭanā) — an event}

\label{lesson:SA-C310-ghatana}



\begin{quote}
Before the new one: say the Sanskrit for main, then the Sanskrit for a news reporter.
\end{quote}

\subsection*{You'll want to know: \sk{घटना}}
\textbf{\sk{घटना}} — \emph{ghaṭanā} — \textquotedblleft{}an event\textquotedblright{}.

\begin{quote}
\sk{अद्य} \sk{नगरे} \sk{एका} \sk{घटना} \sk{घटिता।}
\end{quote}

\begin{quote}
\emph{adya nagare ekā ghaṭanā ghaṭitā.}
\end{quote}

\begin{quote}
Something happened in the city today.
\end{quote}

A feminine noun.

\begin{cousinweb}
\textbf{\sk{घटना}} (\emph{ghaṭanā}) comes from the root √ghaṭ, which also gives \emph{ghaṭate}, \textquotedblleft{}happens\textquotedblright{}; MW gives \emph{ghaṭanā} as exertion, motion and a manner of acting, and an event is its modern sense (MW).
\end{cousinweb}

\begin{grammarlens}[title={What you've built}]
One more word for telling what happened.
\end{grammarlens}

\subsection*{Your turn}
\begin{itemize}
\raggedright
  \item \emph{Say it:} \emph{ghaṭanā}
  \item \emph{Say it:} \emph{ghaṭanā}, once more
  \item \emph{Say it:} \emph{adya nagare ekā ghaṭanā ghaṭitā.} — Something happened in the city today.
  \item \emph{Recall:} say the Sanskrit for main, then the Sanskrit for a news reporter, then say \emph{ghaṭanā} again
\end{itemize}

\begin{tcolorbox}[breakable,colback=teal!4,colframe=teal!35!black,title={Before you move on}]
What does \sk{घटना} mean? (\textquotedblleft{}An event\textquotedblright{}.) Say it once more.
\end{tcolorbox}
\section[\texorpdfstring{\sk{आपतति} (āpatati)}{आपतति (āpatati)}]{\sk{आपतति} (āpatati) — befalls}

\label{lesson:SA-C310-apatati}



\begin{quote}
Before the new one: say the Sanskrit for a news reporter, then the Sanskrit for an event.
\end{quote}

\subsection*{You'll want to know: \sk{आपतति}}
\textbf{\sk{आपतति}} — \emph{āpatati} — \textquotedblleft{}befalls\textquotedblright{}.

\begin{quote}
\sk{ग्रामे} \sk{सहसा} \sk{संकटम्} \sk{आपतति।}
\end{quote}

\begin{quote}
\emph{grāme sahasā saṅkaṭam āpatati.}
\end{quote}

\begin{quote}
Danger suddenly befalls the village.
\end{quote}

Past (it …): \sk{संकटं} \sk{सहसा} \sk{आपतितवत्}, \emph{saṅkaṭaṃ sahasā āpatitavat}.

\begin{cousinweb}
\textbf{\sk{आपतति}} (\emph{āpatati}) is built from \emph{ā}, \textquotedblleft{}towards\textquotedblright{}, and the root √pat, \textquotedblleft{}to fall, fly\textquotedblright{}, as in \emph{patati}; MW gives \emph{ā-pat} as to fall upon, happen, befall (MW).
\end{cousinweb}

\begin{grammarlens}[title={What you've built}]
One more word for telling what happened.
\end{grammarlens}

\subsection*{Your turn}
\begin{itemize}
\raggedright
  \item \emph{Say it:} \emph{āpatati}
  \item \emph{Say it:} \emph{āpatati}, once more
  \item \emph{Say it:} \emph{grāme sahasā saṅkaṭam āpatati.} — Danger suddenly befalls the village.
  \item \emph{Recall:} say the Sanskrit for a news reporter, then the Sanskrit for an event, then say \emph{āpatati} again
\end{itemize}

\begin{tcolorbox}[breakable,colback=teal!4,colframe=teal!35!black,title={Before you move on}]
What does \sk{आपतति} mean? (\textquotedblleft{}Befalls\textquotedblright{}.) Say it once more.
\end{tcolorbox}
\section[\texorpdfstring{\sk{विवरणम्} (vivaraṇam)}{विवरणम् (vivaraṇam)}]{\sk{विवरणम्} (vivaraṇam) — a detailed account}

\label{lesson:SA-C310-vivaranam}



\begin{quote}
Before the new one: say the Sanskrit for an event, then the Sanskrit for befalls.
\end{quote}

\subsection*{You'll want to know: \sk{विवरणम्}}
\textbf{\sk{विवरणम्}} — \emph{vivaraṇam} — \textquotedblleft{}a detailed account\textquotedblright{}.

\begin{quote}
\sk{वार्ताहरः} \sk{घटनायाः} \sk{विवरणं} \sk{ददाति।}
\end{quote}

\begin{quote}
\emph{vārtāharaḥ ghaṭanāyāḥ vivaraṇaṃ dadāti.}
\end{quote}

\begin{quote}
The reporter gives a detailed account of the event.
\end{quote}

A neuter noun.

\begin{cousinweb}
\textbf{\sk{विवरणम्}} (\emph{vivaraṇam}) is built from \emph{vi}, \textquotedblleft{}apart\textquotedblright{}, and the root √vṛ, \textquotedblleft{}to cover\textquotedblright{}: an uncovering, a laying open, and MW also gives \emph{vivaraṇa} as an explanation or exposition; a detailed account of an event is its modern use (MW).
\end{cousinweb}

\begin{grammarlens}[title={What you've built}]
One more word for telling what happened.
\end{grammarlens}

\subsection*{Your turn}
\begin{itemize}
\raggedright
  \item \emph{Say it:} \emph{vivaraṇam}
  \item \emph{Say it:} \emph{vivaraṇam}, once more
  \item \emph{Say it:} \emph{vārtāharaḥ ghaṭanāyāḥ vivaraṇaṃ dadāti.} — The reporter gives a detailed account of the event.
  \item \emph{Recall:} say the Sanskrit for an event, then the Sanskrit for befalls, then say \emph{vivaraṇam} again
\end{itemize}

\begin{tcolorbox}[breakable,colback=teal!4,colframe=teal!35!black,title={Before you move on}]
What does \sk{विवरणम्} mean? (\textquotedblleft{}A detailed account\textquotedblright{}.) Say it once more.
\end{tcolorbox}
\section[\texorpdfstring{\sk{प्रकाशयति} (prakāśayati)}{प्रकाशयति (prakāśayati)}]{\sk{प्रकाशयति} (prakāśayati) — publishes; makes known}

\label{lesson:SA-C310-prakasayati}



\begin{quote}
Before the new one: say the Sanskrit for befalls, then the Sanskrit for a detailed account.
\end{quote}

\subsection*{You'll want to know: \sk{प्रकाशयति}}
\textbf{\sk{प्रकाशयति}} — \emph{prakāśayati} — \textquotedblleft{}publishes; makes known\textquotedblright{}.

\begin{quote}
\sk{पत्रिका} \sk{घटनायाः} \sk{विवरणं} \sk{प्रकाशयति।}
\end{quote}

\begin{quote}
\emph{patrikā ghaṭanāyāḥ vivaraṇaṃ prakāśayati.}
\end{quote}

\begin{quote}
The newspaper publishes a detailed account of the event.
\end{quote}

Past (it …): \sk{पत्रिका} \sk{विवरणं} \sk{प्रकाशितवती}, \emph{patrikā vivaraṇaṃ prakāśitavatī}.

\begin{cousinweb}
\textbf{\sk{प्रकाशयति}} (\emph{prakāśayati}) is built from \emph{pra} and the causative of the root √kāś, \textquotedblleft{}to shine\textquotedblright{}, as in \emph{prakāśa}, \textquotedblleft{}light\textquotedblright{}; MW gives it as to make visible, display, reveal, proclaim, and to publish is its modern use (MW).
\end{cousinweb}

\begin{grammarlens}[title={What you've built}]
One more word for telling what happened.
\end{grammarlens}

\subsection*{Your turn}
\begin{itemize}
\raggedright
  \item \emph{Say it:} \emph{prakāśayati}
  \item \emph{Say it:} \emph{prakāśayati}, once more
  \item \emph{Say it:} \emph{patrikā ghaṭanāyāḥ vivaraṇaṃ prakāśayati.} — The newspaper publishes a detailed account of the event.
  \item \emph{Recall:} say the Sanskrit for befalls, then the Sanskrit for a detailed account, then say \emph{prakāśayati} again
\end{itemize}

\begin{tcolorbox}[breakable,colback=teal!4,colframe=teal!35!black,title={Before you move on}]
What does \sk{प्रकाशयति} mean? (\textquotedblleft{}Publishes; makes known\textquotedblright{}.) Say it once more.
\end{tcolorbox}
\section[\texorpdfstring{\sk{विशेषः} (viśeṣaḥ)}{विशेषः (viśeṣaḥ)}]{\sk{विशेषः} (viśeṣaḥ) — a detail; a particular}

\label{lesson:SA-C310-visesah}



\begin{quote}
Before the new one: say the Sanskrit for a detailed account, then the Sanskrit for publishes.
\end{quote}

\subsection*{You'll want to know: \sk{विशेषः}}
\textbf{\sk{विशेषः}} — \emph{viśeṣaḥ} — \textquotedblleft{}a detail; a particular\textquotedblright{}.

\begin{quote}
\sk{वार्तायां} \sk{कः} \sk{विशेषः} \sk{अस्ति}?
\end{quote}

\begin{quote}
\emph{vārtāyāṃ kaḥ viśeṣaḥ asti?}
\end{quote}

\begin{quote}
What is the special point in the news?
\end{quote}

A masculine noun.

\begin{cousinweb}
\textbf{\sk{विशेषः}} (\emph{viśeṣaḥ}) is built from \emph{vi} and the root √śiṣ, \textquotedblleft{}to leave, distinguish\textquotedblright{}; MW gives \emph{viśeṣa} as a distinction, a peculiar mark, a peculiarity (MW).
\end{cousinweb}

\begin{grammarlens}[title={What you've built}]
One more word for telling what happened.
\end{grammarlens}

\subsection*{Your turn}
\begin{itemize}
\raggedright
  \item \emph{Say it:} \emph{viśeṣaḥ}
  \item \emph{Say it:} \emph{viśeṣaḥ}, once more
  \item \emph{Say it:} \emph{vārtāyāṃ kaḥ viśeṣaḥ asti?} — What is the special point in the news?
  \item \emph{Recall:} say the Sanskrit for a detailed account, then the Sanskrit for publishes, then say \emph{viśeṣaḥ} again
\end{itemize}

\begin{tcolorbox}[breakable,colback=teal!4,colframe=teal!35!black,title={Before you move on}]
What does \sk{विशेषः} mean? (\textquotedblleft{}A detail; a particular\textquotedblright{}.) Say it once more.
\end{tcolorbox}
